\input{../tex-inputs/latex-leseansicht-vorspann}

\section[Arthur und Olga Schnitzler an Gustav Schwarzkopf, 18. 7. 1911]{L04389 Arthur und Olga Schnitzler an Gustav Schwarzkopf, 18. 7. 1911}
\nopagebreak\mylabel{L04389v}
\rehead{ }\normalsize\beginnumbering\briefempfaengerindex{Schwarzkopf, Gustav@\textsc{Schwarzkopf, Gustav}!zzzSchnitzler, Olga@\emph{von Olga Schnitzler}!1911-07-181@{18. 7. 1911}|(be}\briefempfaengerindex{Schwarzkopf, Gustav@\textsc{Schwarzkopf, Gustav}!zzzSchnitzler, Arthur@\emph{von Arthur Schnitzler}!1911-07-181@{18. 7. 1911}|(be}
\toendnotes[C]{\smallbreak\pagebreak[2]}
\correspDesc{Versand  durch Arthur Schnitzler, Olga Schnitzler am 18. 7. 1911 in Sternwartestraße 71
\newline{}Übermittlung  am 18. 7. 1911 in 1/18 Wien
\newline{}Erhalt  durch Gustav Schwarzkopf im Zeitraum [19. 7. 1911 – 22. 7. 1911?] in Gasthof Josef Kloßner}\toendnotes[C]{\smallbreak}
\Standort{DLA, A:Schnitzler, HS.1985.1.1897.}
\physDesc{Bildpostkarte, 286 Zeichen
\newline{}Handschrift Arthur Schnitzler: schwarze Tinte, deutsche Kurrent
\newline{}Handschrift Olga Schnitzler: schwarze Tinte, lateinische Kurrent
\newline{}Versand: Stempel: »\nobreak{}\oindex{XVIII., Währing@\textbf{XVIII., Währing}|pwk}1/1\textcolor{gray}{8} Wien, 1\textcolor{gray}{8}. VII. 11, \textcolor{gray}{7}\nobreak{}«.  }\toendnotes[C]{\smallbreak}\pstart{}{\pb}Herrn Gustav Schwarzkopf\pend{}\pstart{}\textsc{Mautern\oindex{Mautern in Steiermark@\textbf{Mautern in Steiermark}|pw}}\pend{}\pstart{}in \textsc{Steiermark\oindex{Steiermark@\textbf{Steiermark}|pw}}.\pend{}\pstart{}Gaſthof \textsc{Klossner}\oindex{Gasthof Josef Kloßner@\textbf{Gasthof Josef Kloßner}|pw}\pend{}{\bigskip}
\pstart
           \centering{}{\pb}\textcolor{gray}{\textbf{Türkenschanz-Park\oindex{Wien@\textbf{Wien}!XVIII., Währing@\textbf{XVIII., Währing}!Türkenschanzpark@\textbf{Türkenschanzpark}|pw} gegen Gersthof\oindex{Wien@\textbf{Wien}!XVIII., Währing@\textbf{XVIII., Währing}!Gersthof@\textbf{Gersthof}|pw}.}}\hspace*{1.5em}\textcolor{gray}{\textbf{Wien XVIII/I\oindex{XVIII., Währing@\textbf{XVIII., Währing}|pw}.}}\pend
           \vspace{1em}
\pstart
           {\pb}18. 7. 911\pend
           \vspace{0.5em}
\pstart
           herzlicher Gruſs! Wir{ }ſind noch
      hier, ich hole \label{K_L04389-1v}\edtext{am 30. 7.}{\lemma{\textnormal{\emph{am 30. 7.}}}\Cendnote{\textnormal{Vgl. A. S.: \emph{Tagebuch}, 31. 7. 1911.}}}\label{K_L04389-1}{ }Mama\pwindex{Schnitzler, Louise 8.\,7.\,1840 Kőszeg – 9.\,9.\,1911 Wien@\textsc{Schnitzler, Louise} (8.\,7.\,1840 Kőszeg – 9.\,9.\,1911 Wien)|pwv} 
               aus I{\geminationn}sbruck\oindex{Innsbruck@\textbf{Innsbruck}|pw} ab – nach
      Baden\oindex{Baden bei Wien@\textbf{Baden bei Wien}|pw} oder \textsc{Aussee\oindex{Bad Aussee@\textbf{Bad Aussee}|pw}} – eventuell
      gehn wir alle auf den Se{\geminationm}ering\oindex{Semmering@\textbf{Semmering}|pw}.
      Vorläufg nichts{ }ſichres. Aber ich
      find es hier, in Wien\oindex{Wien@\textbf{Wien}|pw},{ }ſehr{ }ſchön.\pend
           
\pstart
           \spacefill\mbox{Ihr A.}{\\[\baselineskip]}\spacefill\mbox{{[}hs. O. Schnitzler:{]} Olga}\pend
           \leftskip=0em{}\selectlanguage{ngerman}\endnumbering\briefempfaengerindex{Schwarzkopf, Gustav@\textsc{Schwarzkopf, Gustav}!zzzSchnitzler, Olga@\emph{von Olga Schnitzler}!1911-07-181@{18. 7. 1911}|)be}\briefempfaengerindex{Schwarzkopf, Gustav@\textsc{Schwarzkopf, Gustav}!zzzSchnitzler, Arthur@\emph{von Arthur Schnitzler}!1911-07-181@{18. 7. 1911}|)be}\mylabel{L04389h}
\begin{anhang}
\end{anhang}\newcommand{\dateiname}{L04389}\newcommand{\titel}{Arthur und Olga Schnitzler an Gustav Schwarzkopf, 18. 7. 1911}\newcommand{\editorInnen}{Herausgegeben von Jahnke, SelmaMüller, Martin Anton}\input{../tex-inputs/latex-leseansicht-abspann}
